% ==============================================================================
% Example 06: Trigger - Trigger Volumes and Overlap Detection
% ==============================================================================

\section{Example 06: Trigger}

\subsection{Overview and Basic Information}

\subsubsection{File Information}

\begin{table}[h]
\centering
\begin{tabular}{|l|l|}
\hline
\textbf{Aspect} & \textbf{Details} \\
\hline
\textbf{Original PhysX Snippet} & \texttt{physx/snippets/snippettriggers/} \\
 & \texttt{SnippetTriggers.cpp} \\
\hline
Original File Size & 450+ lines (3 implementations) \\
\hline
\textbf{Ported PhysXWrapper Example} & \texttt{PhysXWrapper/examples/} \\
 & \texttt{example\_trigger.cpp} \\
\hline
Ported File Size & 486 lines \\
\hline
\textbf{Core Wrapper Class} & \texttt{RigidBody/RigidBodyTrigger.h/cpp} \\
\hline
\textbf{Difficulty Level} & ⭐⭐ (Intermediate) \\
\hline
\textbf{Functional Module} & RigidBody - Trigger System \\
\hline
\textbf{Primary Features} & Trigger volumes, Enter/exit events, CCD triggers \\
\hline
\end{tabular}
\caption{Trigger Example Basic Information}
\end{table}

\subsubsection{Dependencies}

\textbf{Original PhysX Dependencies:}
\begin{itemize}
    \item \texttt{PxShapeFlag::eTRIGGER\_SHAPE} - Native trigger flag
    \item \texttt{PxSimulationEventCallback} - Trigger event callback
    \item \texttt{PxPairFlag::eTRIGGER\_DEFAULT} - Trigger pair flags
    \item Filter shader/callback for trigger emulation
\end{itemize}

\textbf{PhysXWrapper Dependencies:}
\begin{itemize}
    \item \texttt{Core/PhysXCore.h} - Core initialization
    \item \texttt{RigidBody/RigidBodyTrigger.h} - Trigger manager
\end{itemize}

\subsubsection{Learning Objectives}

\begin{enumerate}
    \item \textbf{Trigger Volumes:} Non-physical overlap detection zones
    \item \textbf{Event Callbacks:} Enter/exit event handling
    \item \textbf{CCD Integration:} Triggers with fast-moving objects
    \item \textbf{Multiple Implementations:} Native vs filter-based triggers
    \item \textbf{Trigger-Trigger Detection:} Detecting trigger overlaps
\end{enumerate}

%------------------------------------------------------------------------------

\subsection{Functional Description}

\subsubsection{What are Triggers?}

\textbf{Triggers} are special collision shapes that detect overlaps without generating physical collision responses. They are used for:

\begin{itemize}
    \item \textbf{Zone Detection:} Detecting objects entering/exiting areas
    \item \textbf{Checkpoints:} Racing game checkpoint systems
    \item \textbf{Proximity Sensors:} AI awareness zones
    \item \textbf{Pickup Items:} Collectible item detection
\end{itemize}

\textbf{Key Properties:}
\begin{itemize}
    \item No collision response (objects pass through)
    \item Generate onTriggerEnter/onTriggerExit events
    \item Compatible with all shape types (box, sphere, capsule, mesh)
    \item Can be attached to static or dynamic actors
\end{itemize}

\subsubsection{Physics Concepts}

\paragraph{Overlap Detection}

Triggers use the same broad-phase and narrow-phase collision detection as regular shapes, but skip the constraint solver:

\begin{equation}
\text{Trigger Path:} \quad \text{BroadPhase} \rightarrow \text{NarrowPhase} \rightarrow \text{Event Callback}
\end{equation}

\begin{equation}
\text{Normal Path:} \quad \text{BroadPhase} \rightarrow \text{NarrowPhase} \rightarrow \text{Solver} \rightarrow \text{Integration}
\end{equation}

\paragraph{CCD for Triggers}

Standard triggers can miss fast-moving objects due to discrete collision detection. With CCD enabled:

\begin{equation}
\text{CCD Trigger:} \quad \text{Swept AABB} \rightarrow \text{Conservative Advancement} \rightarrow \text{Event}
\end{equation}

This ensures even thin triggers (e.g., 0.01m thick) detect fast bullets.

%------------------------------------------------------------------------------

\subsection{Code Comparison}

\subsubsection{Architecture Differences}

\textbf{Original PhysX:}
\begin{itemize}
    \item Three separate implementations (REAL\_TRIGGERS, FILTER\_SHADER, FILTER\_CALLBACK)
    \item Global scenario selection
    \item Manual filter shader/callback setup
    \item Direct PxSimulationEventCallback inheritance
\end{itemize}

\textbf{PhysXWrapper:}
\begin{itemize}
    \item Unified \texttt{RigidBodyTrigger} manager
    \item Configuration via \texttt{TriggerConfig} enum
    \item Simplified callback interface using std::function
    \item Built-in shape creation helpers
\end{itemize}

\subsubsection{Side-by-Side Comparison}

\textbf{Original: Manual Filter Shader Setup}
\begin{lstlisting}[caption={Original PhysX filter shader for triggers}]
static PxFilterFlags triggersUsingFilterShader(
    PxFilterObjectAttributes attributes0, PxFilterData filterData0,
    PxFilterObjectAttributes attributes1, PxFilterData filterData1,
    PxPairFlags& pairFlags,
    const void* constantBlock, PxU32 constantBlockSize)
{
    // Detect trigger using filter data
    const bool isTriggerPair = isTrigger(filterData0) ||
                                isTrigger(filterData1);

    if(isTriggerPair) {
        pairFlags = PxPairFlag::eTRIGGER_DEFAULT;
        if(usesCCD())
            pairFlags |= PxPairFlag::eDETECT_CCD_CONTACT;
        return PxFilterFlag::eDEFAULT;
    } else {
        pairFlags = PxPairFlag::eCONTACT_DEFAULT;
        return PxFilterFlag::eDEFAULT;
    }
}

// Scene setup
PxSceneDesc sceneDesc(physics->getTolerancesScale());
sceneDesc.gravity = PxVec3(0.0f, -9.81f, 0.0f);
sceneDesc.cpuDispatcher = dispatcher;
sceneDesc.filterShader = triggersUsingFilterShader;
PxScene* scene = physics->createScene(sceneDesc);
\end{lstlisting}

\textbf{Wrapper: Simplified Configuration}
\begin{lstlisting}[caption={Wrapper: Config-driven trigger setup}]
// Create trigger manager
RigidBodyTrigger triggerManager;

// Set callback
triggerManager.setTriggerCallback([](const TriggerEvent& event) {
    if (event.type == TriggerEventType::ENTER) {
        std::cout << "Object entered trigger!" << std::endl;
    } else {
        std::cout << "Object exited trigger!" << std::endl;
    }
});

// Configure trigger behavior
TriggerConfig config;
config.implementation = TriggerImplementation::FILTER_SHADER;
config.enableCCD = true;
config.detectTriggerTrigger = false;

// Create scene with automatic filter setup
PxTolerancesScale scale;
PxSceneDesc sceneDesc = triggerManager.createSceneDesc(scale, config);
sceneDesc.cpuDispatcher = PxDefaultCpuDispatcherCreate(2);
PxScene* scene = physics.getPhysics()->createScene(sceneDesc);

// Create trigger shape easily
PxRigidStatic* triggerActor = triggerManager.createBoxTrigger(
    physics.getPhysics(),
    scene,
    material,
    PxVec3(0, 10, 0),     // Position
    PxVec3(10, 1, 10),    // Half extents
    config
);
\end{lstlisting}

%------------------------------------------------------------------------------

\subsection{Core Class Analysis}

\subsubsection{TriggerConfig Structure}

\begin{lstlisting}[caption={Trigger configuration}]
enum class TriggerImplementation {
    NATIVE,          // PxShapeFlag::eTRIGGER_SHAPE (default)
    FILTER_SHADER,   // Emulated via filter shader (supports CCD)
    FILTER_CALLBACK  // Emulated via callback (custom logic)
};

struct TriggerConfig {
    TriggerImplementation implementation = TriggerImplementation::NATIVE;
    bool enableCCD = false;               // CCD for fast objects
    bool detectTriggerTrigger = false;    // Trigger-trigger overlaps
};
\end{lstlisting}

\subsubsection{TriggerEvent Structure}

\begin{lstlisting}[caption={Trigger event data}]
enum class TriggerEventType {
    ENTER,  // Object entered trigger
    EXIT    // Object exited trigger
};

struct TriggerEvent {
    TriggerEventType type;
    PxShape* triggerShape;    // The trigger shape
    PxShape* otherShape;      // The shape that triggered it
    PxActor* triggerActor;
    PxActor* otherActor;
};
\end{lstlisting}

\subsubsection{Key Methods}

\begin{lstlisting}[caption={RigidBodyTrigger core methods}]
class RigidBodyTrigger {
public:
    // Set trigger callback
    void setTriggerCallback(
        std::function<void(const TriggerEvent&)> callback
    );

    // Create scene descriptor with trigger support
    PxSceneDesc createSceneDesc(
        const PxTolerancesScale& scale,
        const TriggerConfig& config
    );

    // Create trigger shapes
    PxRigidStatic* createBoxTrigger(
        PxPhysics* physics, PxScene* scene, PxMaterial* material,
        const PxVec3& position, const PxVec3& halfExtents,
        const TriggerConfig& config
    );

    PxRigidStatic* createSphereTrigger(
        PxPhysics* physics, PxScene* scene, PxMaterial* material,
        const PxVec3& position, PxReal radius,
        const TriggerConfig& config
    );

    PxRigidStatic* createCapsuleTrigger(
        PxPhysics* physics, PxScene* scene, PxMaterial* material,
        const PxVec3& position, PxReal radius, PxReal halfHeight,
        const TriggerConfig& config
    );

    // Create standalone trigger shape
    PxShape* createTriggerShape(
        PxPhysics* physics, const PxGeometry& geometry,
        PxMaterial* material, const TriggerConfig& config,
        bool isExclusive = true
    );

    // Utility
    static bool isTriggerShape(
        PxShape* shape, const TriggerConfig& config
    );
};
\end{lstlisting}

%------------------------------------------------------------------------------

\subsection{Visual Diagrams}

\subsubsection{UML Class Diagram}

\begin{figure}[h]
\centering
\begin{tikzpicture}[
    node distance=1.5cm,
    class/.style={rectangle, draw=black, thick, minimum width=4cm, minimum height=1cm},
    struct/.style={rectangle, draw=blue!70, thick, minimum width=3cm}
]

\node[class] (trigger) {
    \begin{tabular}{c}
    \textbf{RigidBodyTrigger} \\
    \hline
    - m\_callback: function \\
    - m\_eventCallback: EventCallback* \\
    \hline
    + setTriggerCallback(...) \\
    + createSceneDesc(...): PxSceneDesc \\
    + createBoxTrigger(...): Actor* \\
    + createSphereTrigger(...): Actor* \\
    + createTriggerShape(...): Shape* \\
    + isTriggerShape(...): bool \\
    \end{tabular}
};

\node[struct, below left=0.8cm and 0.5cm of trigger] (config) {
    \begin{tabular}{c}
    \textbf{TriggerConfig} \\
    \hline
    + implementation: Enum \\
    + enableCCD: bool \\
    + detectTriggerTrigger: bool
    \end{tabular}
};

\node[struct, below right=0.8cm and 0.5cm of trigger] (event) {
    \begin{tabular}{c}
    \textbf{TriggerEvent} \\
    \hline
    + type: EventType \\
    + triggerShape: PxShape* \\
    + otherShape: PxShape* \\
    + triggerActor: PxActor* \\
    + otherActor: PxActor*
    \end{tabular}
};

\draw[->, dashed] (trigger) -- (config) node[midway, left] {uses};
\draw[->, dashed] (trigger) -- (event) node[midway, right] {emits};

\end{tikzpicture}
\caption{RigidBodyTrigger UML Diagram}
\end{figure}

\subsubsection{Trigger Event Flow}

\begin{figure}[h]
\centering
\begin{tikzpicture}[
    node distance=1cm,
    process/.style={rectangle, rounded corners, draw=blue!70, fill=blue!10, thick, minimum width=2.5cm, minimum height=0.7cm},
    decision/.style={diamond, draw=orange!70, fill=orange!10, thick, minimum width=1.8cm}
]

\node[process] (sim) {Physics Simulation};
\node[process, below=of sim] (broad) {Broad Phase};
\node[process, below=of broad] (narrow) {Narrow Phase};
\node[decision, below=of narrow] (overlap) {Overlap?};
\node[process, below left=1cm and 0.5cm of overlap] (enter) {Fire ENTER Event};
\node[process, below right=1cm and 0.5cm of overlap] (exit) {Fire EXIT Event};
\node[process, below=2cm of overlap] (callback) {User Callback};

\draw[->, thick] (sim) -- (broad);
\draw[->, thick] (broad) -- (narrow);
\draw[->, thick] (narrow) -- (overlap);
\draw[->, thick] (overlap) -- node[left] {New} (enter);
\draw[->, thick] (overlap) -- node[right] {Lost} (exit);
\draw[->, thick] (enter) -- (callback);
\draw[->, thick] (exit) -- (callback);

\end{tikzpicture}
\caption{Trigger Event Flow}
\end{figure}

%------------------------------------------------------------------------------

\subsection{Usage Methodology}

\subsubsection{Basic Trigger Usage}

\begin{lstlisting}[caption={Simple trigger zone}]
RigidBodyTrigger triggerManager;

// Set callback
triggerManager.setTriggerCallback([](const TriggerEvent& event) {
    if (event.type == TriggerEventType::ENTER) {
        std::cout << "Checkpoint reached!" << std::endl;
    }
});

// Create trigger
TriggerConfig config;
config.implementation = TriggerImplementation::NATIVE;

PxRigidStatic* checkpoint = triggerManager.createBoxTrigger(
    physics, scene, material,
    PxVec3(0, 1, 50),      // Checkpoint position
    PxVec3(10, 5, 2),      // Size
    config
);
\end{lstlisting}

\subsubsection{Triggers with CCD}

\begin{lstlisting}[caption={CCD trigger for fast objects}]
TriggerConfig config;
config.implementation = TriggerImplementation::FILTER_SHADER;
config.enableCCD = true;  // Enable CCD

// Create very thin trigger (would be missed without CCD)
PxRigidStatic* sensor = triggerManager.createBoxTrigger(
    physics, scene, material,
    PxVec3(0, 10, 0),
    PxVec3(20, 0.01f, 20),  // 0.01m thin!
    config
);

// Create fast-moving bullet
PxRigidDynamic* bullet = physics.createRigidDynamic(
    PxTransform(0, 20, 0));
bullet->setRigidBodyFlag(PxRigidBodyFlag::eENABLE_CCD, true);
bullet->setLinearVelocity(PxVec3(0, -150, 0));  // Very fast
\end{lstlisting}

%------------------------------------------------------------------------------

\subsection{Feature Completeness}

\begin{table}[h]
\centering
\begin{tabular}{|l|c|c|}
\hline
\textbf{Feature} & \textbf{Original} & \textbf{Wrapper} \\
\hline
Native triggers & ✓ & ✓ \\
Filter shader triggers & ✓ & ✓ \\
Filter callback triggers & ✓ & ✓ \\
CCD support & ✓ & ✓ \\
Trigger-trigger detection & ✓ & ✓ \\
Unified API & ✗ & ✓ \\
Config-driven & ✗ & ✓ \\
Shape creation helpers & ✗ & ✓ \\
Lambda callbacks & ✗ & ✓ \\
\hline
\end{tabular}
\caption{Feature Comparison}
\end{table}

%------------------------------------------------------------------------------

\subsection{Quality Assessment}

\begin{table}[h]
\centering
\begin{tabular}{|l|c|l|}
\hline
\textbf{Category} & \textbf{Rating} & \textbf{Justification} \\
\hline
Functional Fidelity & ⭐⭐⭐⭐⭐ & All 3 implementations supported \\
Code Quality & ⭐⭐⭐⭐⭐ & Clean unified API \\
Usability & ⭐⭐⭐⭐⭐ & Simplified configuration \\
Documentation & ⭐⭐⭐⭐⭐ & Clear examples \\
\hline
\end{tabular}
\caption{Quality Assessment}
\end{table}

\textbf{Approval Status:} \textcolor{green}{\textbf{APPROVED}} ✓

%==============================================================================
% End of Example 06: Trigger
%==============================================================================
